\newglossaryentry{computer}
{
  name=computer,
  description={is a programmable machine that receives input,
               stores and manipulates data, and provides
               output in a useful format}
}

\newacronym[longplural={Frames per Second}]{fpsLabel}{FPS}{Frame per Second}
\newacronym{lvm}{LVM}{Logical Volume Manager}

% ─── GenAI-Flickr Acronyms ──────────────────────────────────────────────

\newacronym{ssm}{SSM}{State Space Model}
\newacronym{smpl}{SMPL}{Skinned Multi-Person Linear Model}
\newacronym{rom}{ROM}{Range of Motion}
\newacronym{rsi}{RSI}{Root State Injection}
\newacronym{lbs}{LBS}{Linear Blend Skinning}
\newacronym{fk}{FK}{Forward Kinematics}
\newacronym{dof}{DoF}{Depth of Field}
\newacronym{vram}{VRAM}{Video Random Access Memory}
\newacronym{gpu}{GPU}{Graphics Processing Unit}
\newacronym{mlp}{MLP}{Multi-Layer Perceptron}
\newacronym{vae}{VAE}{Variational Autoencoder}
\newacronym{rl}{RL}{Reinforcement Learning}
\newacronym{rnn}{RNN}{Recurrent Neural Network}
\newacronym{sd}{SD}{Stable Diffusion}

% ─── GenAI-Flickr Glossary Terms ────────────────────────────────────────

\newglossaryentry{mamba}
{
  name=Mamba,
  description={un model de spațiu de stări selectiv cu complexitate liniară,
               propus de Gu și Dao (2023), utilizat în MotionSSM pentru
               modelarea secvențelor temporale de mișcare}
}

\newglossaryentry{physicsssm}
{
  name=PhysicsSSM,
  description={contribuția originală a acestei lucrări --- un modul SSM care
               integrează constrângeri fizice (gravitație, contact, jerk)
               prin intermediul unei porți sigmoid învățate}
}

\newglossaryentry{kitml}
{
  name=KIT-ML,
  description={KIT Motion-Language --- un dataset de motion capture cu 4.888
               de secvențe la 20 fps, fiecare cu descrieri text în engleză,
               folosit pentru antrenarea și evaluarea mișcării umane}
}

\newglossaryentry{pybullet}
{
  name=PyBullet,
  description={un motor de simulare fizică open-source pentru robotică,
               jocuri și machine learning, utilizat în M5 pentru simularea
               gravitației, coliziunilor și contactelor}
}

\newglossaryentry{faiss}
{
  name=FAISS,
  description={Facebook AI Similarity Search --- o bibliotecă pentru
               căutarea eficientă de similaritate în spații de embedding
               de dimensiune mare, folosită în M1 pentru knowledge retrieval}
}

\newglossaryentry{sentencebert}
{
  name=Sentence-BERT,
  description={un model de embedding bazat pe BERT care produce
               reprezentări vectoriale (384-dim) ale propozițiilor,
               folosit pentru comparația semantică text-mișcare}
}

\newglossaryentry{shapE}
{
  name=Shap-E,
  description={un model generativ de la OpenAI care produce funcții
               implicite 3D condiționate de text, folosit în M3 pentru
               generarea automată de modele 3D}
}

\newglossaryentry{controlnet}
{
  name=ControlNet,
  description={o arhitectură care adaugă control condiționat
               (schelet OpenPose, edge maps) la modele de difuzie
               text-to-image, folosită în M7 pentru randare realistă}
}

\newglossaryentry{humanml3d}
{
  name=HumanML3D,
  description={un format de reprezentare a mișcării umane cu 251 de
               dimensiuni per cadru (rotații articulare, poziție root,
               viteze), standard în cercetarea de motion generation}
}
